\section{Method}
\label{sec:method}

This section describes the pipeline in Fig.~\ref{fig:pipeline}: tile the photo at a common leaf scale, describe each
tile once with a frozen encoder, and combine the descriptions with a small trained attention head.

\begin{figure}[t]
\centering
\begin{tikzpicture}[
  font=\scriptsize,
  box/.style={draw, rounded corners=1.5pt, align=center, minimum height=0.85cm, text width=6.9cm, inner sep=2pt},
  frozen/.style={box, fill=black!6},
  trained/.style={box, fill=black!22},
  outbox/.style={draw, rounded corners=1.5pt, align=center, minimum height=0.8cm, text width=3.1cm, inner sep=2pt},
  arrow/.style={-{Latex[length=1.6mm]}, thick}]
  \node[box] (photo) {Rice leaf photo (734--3{,}120 px per side), photo-level label for training only};
  \node[box, below=0.28cm of photo] (tile) {\textbf{1. Rescale and tile.} Common leaf scale, then 49--100 tiles of $224\times224$ px (no learning)};
  \node[frozen, below=0.28cm of tile] (enc) {\textbf{2. Describe.} Frozen DINOv3 ViT-B/16: 1{,}536 numbers per tile, computed once and cached};
  \node[trained, below=0.28cm of enc] (att) {\textbf{3. Attend and decide.} Gated-attention head ($0.23$ M trained parameters): weight $a_i$ per tile, weighted sum, four class scores};
  \node[outbox, below=0.34cm of att, xshift=-1.85cm] (lab) {Class: healthy, leaf blast, brown spot or hispa};
  \node[outbox, below=0.34cm of att, xshift=1.85cm] (map) {Heat map: weights $a_i$ on the tile grid};
  \draw[arrow] (photo) -- (tile);
  \draw[arrow] (tile) -- node[right, xshift=1pt] {tiles} (enc);
  \draw[arrow] (enc) -- node[right, xshift=1pt] {descriptions} (att);
  \draw[arrow] (att.south) ++(-1.85,0) -- (lab.north);
  \draw[arrow] (att.south) ++(1.85,0) -- (map.north);
\end{tikzpicture}
\caption{The pipeline. White boxes: no learning; light grey: frozen encoder; dark grey: the only trained part.
Only the photo-level label is used for training.}
\label{fig:pipeline}
\end{figure}

\subsection{Common leaf scale and tiles}
A 224-pixel tile covers a different amount of leaf in photos of 734 and 3,120 pixels, so we first brought every photo to a
common leaf scale. The scale rule is described in Section~\ref{sec:setup}; it predicts the leaf width
$\hat w(s)=\exp(b+a\log s)$ in pixels from the shorter side $s$, with $a=0.776456$ and $b=-0.557512$, and rescales the photo
by $194.9/\hat w(s)$, where 194.9 px is the median measured leaf width (Section~\ref{sec:setup:data}). The rescaled shorter sides run from 1,488 to
2,056 px. For a side of length $L$ we cut $n=\lceil L/224\rceil$ square $224$-px tiles per axis, with starts evenly spaced
from $0$ to $L-224$ (integer pixel positions by floor division), so the tiles cover the whole photo, background
included. This gives 49--100 tiles per photo and 265{,}349 in all. We call the result ``high-detail tiles at a common
leaf scale''; tiles at the stored resolution (the \emph{native scale}) and tiles of Otsu-cleaned photos are ablations
(Section~\ref{sec:results}).

\subsection{Frozen description}
A frozen DINOv3 ViT-B/16 encoder \cite{simeoni2025dinov3} turned each tile into tokens: one class token, four register
tokens \cite{darcet2024registers} and the patch tokens. The descriptor of a tile is the 768-number class token $\mathbf c_i$
joined with the 768-number mean $\mathbf m_i$ of the patch tokens, which excludes the class and register tokens
($\mathbf h_i=[\mathbf c_i;\mathbf m_i]$, 1{,}536 numbers; descriptors using only $\mathbf c_i$ or only $\mathbf m_i$ are
ablations). Descriptions were computed once and cached for eight versions of each photo (the clean photo and augmented versions;
only the clean one is used for validation and test); the encoder (85.7 M parameters) is never updated, so it cannot memorise the photos.

\subsection{Gated-attention head}
For a bag of descriptors $\{\mathbf h_i\}$ a learned layer gives
$\mathbf g_i=\operatorname{Dropout}(\operatorname{ReLU}(W_p\mathbf h_i+\mathbf b_p))$. The gated attention of Ilse
\textit{et al.}\ \cite{ilse2018} scores and normalises the tiles and forms the photo representation $\mathbf z$:
\begin{align}
 e_i &= \mathbf v^{\top}\bigl[\tanh(V\mathbf g_i+\mathbf b_V)\odot\sigma(U\mathbf g_i+\mathbf b_U)\bigr],\\
 a_i &= \frac{\exp(e_i)}{\sum_j\exp(e_j)},\qquad
 \mathbf z=\sum_i a_i\mathbf g_i .
\end{align}
A linear layer maps $\mathbf z$ to four class scores. Only the projection, the attention parameters and this layer are
trained (0.23 M parameters). During training we dropped tiles at random with a tuned probability and always kept at
least one. The weights $a_i$ sum to one over the photo; painted on the tile footprints (overlaps averaged, scaled to
each photo's maximum), they form the heat map. They describe where the head looked and are not a segmentation. We call the full model Attention MIL.

\subsection{Variants}
Replacing the attention with the mean or the elementwise maximum of the $\mathbf g_i$ gave the pooling controls
(Mean pooling over tiles, Max pooling over tiles). Applying the head to one description of the whole photo, at 224 or
512 px, gave the whole-photo controls (Same head, whole photo); a bag of one tile has no attention to learn. The
Committee averages the class probabilities of ConvNeXt-T (448 px) and Attention MIL, and the Small model (16 tiles) keeps only the 16 tiles ranked
highest by a simple colour rule, to cut the number of encoder passes (Section~\ref{sec:results}).
